\section{Recipes}
This handbook is called the Cookbook of Matfyz, but it would be no cookbook if
it did not actually contain any recipes. Below we have a handful of recipes
prepared for you.
The preparation of these meals is very easy and fast so you can cook them any
time, e.g. when you will be hungry after a long day at school.


\subsubsection*{Sausage goulash}
To prepare three or four serves we will need 1 bigger onion, 200 g sausages
(wieners or \uv{buřty}), 6 potatoes, oil, 4 spoons of flour (\uv{hladká mouka}),
3-4 spoons of paprika (spice), cumin, pepper, bay leaf (3 pcs), salt.

Method: Finely chop the onion, put it in a pot with heated oil and fry until
golden, then dust it with paprika and flour, stir and pour cold! water over it.
Add peeled and sliced potatoes, mix in pepper, bay leaf, cumin and salt and cook
until the potatoes are soft.
Then add diced and peeled sausages, season to taste and cook for another 2
minutes. If the sausage goulash is not thick enough, add flour mixed with water
and let it cook for a while. Serve with bread.


\subsubsection*{Egg bread}
Perfect option for a quick dinner -- and also a convenient way to get rid off
stale bread. You will need bread (stale bread is actually better), eggs (how
many depends on the amount of bread), salt and oil or butter.

Method: Break the eggs into a bowl, mix with salt and a bit of water/milk and
beat everything well (it is best to use a fork to do this).
Then take slices of the bread, dip them in the egg, place over a hot pan and
cook until done. You can either flip the bread or just leave it on one side;
also various seasoning might be added to the egg (pepper, garlic, cheese,
ham...).
Serve it with whatever you like, e.g. ketchup, mustard, onion...


\subsubsection*{Couscous}
Preparing couscous is very easy, just add hot water and wait for a while. You
can eat it either sweet (with cocoa or canned fruit) or salty (with vegetable,
meat, tuna, cheese or whatever left overs you might have). It might even be used
as a side dish with some sauce.


\subsubsection*{Pasta pomodoro}
This recipe is a variation of spaghetti with ketchup. It uses only a few
ingredients and can save you from hunger even if you did not have the time to
buy groceries.

Ingredients: pasta, tomato paste, onion, (garlic), cheese (Parmesan cheese is
the best to go for but any will do), basil/oregano and salt.

Method: Boil the pasta in hot water. Fry the finely chopped onion in another
pan, add garlic, pour tomato paste over it and season to taste.

Note: It tastes better if you make it from fresh tomatoes but storing tomato
paste is much easier as you do not need to store it in the fridge and it stays
fresh longer.


\subsubsection*{Rice with ...}
One of all time top foods that students cook. Cook rice and add anything you
find.
For example you might fry the onion and add frozen vegetables. Season with salt,
pepper and marjoram and you are ready to eat. Or you can eat the rice with
aubergine and bell peppers.
Also the spices might be much different: try chilli, thyme, curry, basil or
mustard. You can also add some meat or sausage and sprinkle cheese over the
meal.
Generally speaking, add any left overs you have and make a nice, satisfying
meal.


\subsubsection*{Tips and tricks}
Following recipes are meant to be used in the cases of emergency (or by truly
\uv{gourmet} cooks). Do you have spaghetti or eggs but not a pot?

\subsubsubsection{Hardboiling eggs in the microwave}
Did you think that hardboiling eggs in the microwave is just a joke? Do you know
someone whose attempt to do this ended badly? Do you want to be better than
them? Lets try out this simple recipe.

Pour a bit of oil into a mug (rather smaller as the eggs are not usually big),
then break an egg in the same cup, season with salt and put it in the microwave.
Microwave for 45 s $\pm$ 10 s depending on the strenght of your microwave and
required consistence. Eat straight from the mug with a spoon.


\subsubsubsection{Cooking spaghetti in a jug kettle}
It is possible to cook spaghetti in a jug cattle, all you need is just a little
bit of patience. Pour water into the kettle and add spaghetti. About half of the
kettle needs to stay empty. Turn on the kettle and once the spaghetti start to
soften, push them into the water until they are all the way under surface. Then
cook until the spaghetti are soft.


\subsubsubsection{Cooking tips}
\begin{itemize}
\item frying - the oil or butter needs to be heated up so anything you are
frying would not soak up that much oil
\item canned meat - convenient ingredient to add to rice, pasta... you can buy a
good, reasonably priced can in Lidl
\item sweet meals (cakes, pancakes) - almost all of them taste better if you add
a pinch of salt
\item tomato sauce - you would not believe how good it tastes with a touch of
cinnamon
\item 1 tablespoon of flour/sugar equals about 25 g
\item meat broth - put the meat into cold water rather than hot, it will release
more juice
\item peeling the tomatoes - it is easier if you pour hot water over the
tomatoes first
\end{itemize}

\subsubsubsection{Refectories}
If your attempts to cook anything fail, you can still visit any of the
refectories on workdays :)

\subsubsubsection{Instructions for cleaning the jug kettle}
You will not get anything to eat following these instructions but you sure will
appreciate the outcome as the water in our dormitory is rich in minerals.
After about a month of using the jug kettle you will notice that a layer of iron
settled on the bottom of it and after half a year it will start to be difficult
to ignore this problem. If you do not want to throw it away and buy a new one,
you will have to clean it.

It is very simple. Pour water and vinegar into the kettle in a ratio of 1:1 and
boil. Let it work for about an hour and then wash out the kettle with clean
water. You might need to repeat this a few times.
If you want to try this with something stronger than vinegar, add acetic acid to
warm water and let it work for a while.